\newcommand{\ModuleID}{EL-A4}
\newcommand{\ModuleTitle}{Rhetoric and Periodic Prose}
\newcommand{\ModuleStage}{Stage IV — Advanced Reading and Composition}
\newcommand{\ModuleUnit}{A4}
\newcommand{\ModuleOutcome}{Prove a rhetorical pattern from exact grammatical members and explain how it orders, delays, balances, or foregrounds thought}
\newcommand{\ModulePrerequisites}{EL-A3 and complete syntactic control of long prayer periods}
\newcommand{\ModuleExtensionPractice}{Worksheets IV.13 and IV.15 remain optional figure-diagnostic and rhetorical-imitation practice}
\newcommand{\ModuleNext}{EL-A5, Textual Comparison and Commentary}
\newcommand{\ModuleMastery}{On an unfamiliar prayer period, account for every clause and non-finite form, prove each proposed figure from corresponding Latin members, and explain its local function rather than merely naming terminology.}
\newcommand{\ModuleDelayNote}{}
\newcommand{\ModuleLessonSource}{\CourseRoot/shared/content/volume4/all}
\newcommand{\ModuleExerciseSource}{\CourseRoot/shared/exercises/volume3_4/volume4}
\newcommand{\ModuleWorkedBank}{IV}
\newcommand{\ModuleExerciseBank}{IV}
\newcommand{\ModuleWorksheetVariants}{B,D}
\newcommand{\ModuleScopeNote}{The packet treats literary function in selected prose and chant; it does not make unsupported claims about quantitative clausulae or authorial intention.}
\newcommand{\ModulePractice}{%
  \SubmoduleExtension{A functional inventory}
    {The same visible repetition may serve balance, accumulation, contrast, or closure depending on its grammatical members.}
    {Choose one figure already printed in the lesson and fill four rows: exact form, syntactic role, corresponding member, and local development.}
    {Remove the figure's name and explain whether the analysis remains intelligible.}
    {A strong model still identifies what repeats or reverses, how members correspond, and what changes across them. If the account collapses without the label, it has not shown function.}
  \SubmoduleExtension{Prose rhythm}
    {A cautious rhythm note begins with clause boundaries, stress or sound observations, and a reproducible comparison rather than a vague claim of solemnity.}
    {On the lesson's period, mark syntactic cola and repeated sounds, then write one claim you can support and one stronger historical or metrical claim you must withhold.}
    {Identify the grammatical boundary that anchors each proposed colon.}
    {The model aligns cola with finite or subordinate units and cites audible or visual recurrence. It withholds quantitative, authorial, or causal conclusions not established by the course evidence.}
  \SubmoduleExtension{Worked period process}
    {A full process moves from verbs to dependencies to rhetorical function.}
    {Using the complete period selected for Worksheet IV.14, produce four layers: verbal inventory, indented clause tree, parallel-member overlay, and three-sentence functional account.}
    {Which layer must be correct before a claim of suspension can pass?}
    {The clause tree must correctly locate the main completion and every dependent before the overlay or functional account is scored. Suspension is proved by the grammatical material delaying that completion.}
}
